\documentclass[11pt]{article}
\usepackage[margin=1in]{geometry}
\usepackage{amsmath,amssymb,amsthm,mathtools}
\usepackage{microtype}
\input{glyphtounicode}
\pdfgentounicode=1
\pdfglyphtounicode{mu}{03BC}
\pdfglyphtounicode{bardbl}{2016}
\pdfglyphtounicode{parenleftbig}{0028}
\pdfglyphtounicode{parenrightbig}{0029}
\pdfglyphtounicode{parenleftBig}{0028}
\pdfglyphtounicode{parenrightBig}{0029}
\pdfglyphtounicode{parenleftbigg}{0028}
\pdfglyphtounicode{parenrightbigg}{0029}
\pdfglyphtounicode{floorleftbigg}{230A}
\pdfglyphtounicode{floorrightbigg}{230B}
\pdfglyphtounicode{braceleftbigg}{007B}
\pdfglyphtounicode{bracerightbigg}{007D}
\pdfglyphtounicode{bracketleftbt}{23A3}
\pdfglyphtounicode{bracketrightbt}{23A6}
\pdfglyphtounicode{bracketleftex}{23A2}
\pdfglyphtounicode{bracketrightex}{23A5}
\pdfglyphtounicode{circleplusdisplay}{2A01}
\pdfglyphtounicode{summationdisplay}{2211}
\pdfglyphtounicode{summationtext}{2211}
\pdfglyphtounicode{productdisplay}{220F}
\pdfglyphtounicode{producttext}{220F}
\pdfglyphtounicode{tildewide}{0303}
\usepackage{xcolor}
\usepackage{listings}
\usepackage{xurl}
\usepackage[colorlinks=true,linkcolor=blue!45!black,citecolor=blue!45!black,urlcolor=blue!45!black]{hyperref}
\usepackage{accsupp}
\NewCommandCopy{\OriginalHookrightarrow}{\hookrightarrow}
\NewCommandCopy{\OriginalMapsto}{\mapsto}
\NewCommandCopy{\OriginalLongmapsto}{\longmapsto}
\NewCommandCopy{\OriginalLStroke}{\l}
\renewcommand{\hookrightarrow}{\mathrel{%
  \BeginAccSupp{method=hex,unicode,ActualText=21AA}%
  \OriginalHookrightarrow\EndAccSupp{}}}
\renewcommand{\mapsto}{\mathrel{%
  \BeginAccSupp{method=hex,unicode,ActualText=21A6}%
  \OriginalMapsto\EndAccSupp{}}}
\renewcommand{\longmapsto}{\mathrel{%
  \BeginAccSupp{method=hex,unicode,ActualText=27FC}%
  \OriginalLongmapsto\EndAccSupp{}}}
\renewcommand{\l}{%
  \BeginAccSupp{method=hex,unicode,ActualText=0142}%
  \OriginalLStroke\EndAccSupp{}}
\hypersetup{pdftitle={Foulkes' conjecture for the sixth symmetric power},pdfauthor={OpenAI}}
\pdfinfoomitdate=1
\pdftrailerid{}
\pdfsuppressptexinfo=15
\newtheorem{theorem}{Theorem}[section]
\newtheorem{lemma}[theorem]{Lemma}
\newtheorem{proposition}[theorem]{Proposition}
\newtheorem{corollary}[theorem]{Corollary}
\newtheorem{verification}[theorem]{Finite verification}
\theoremstyle{definition}
\newtheorem{definition}[theorem]{Definition}
\theoremstyle{remark}
\newtheorem{remark}[theorem]{Remark}
\DeclareMathOperator{\Sym}{Sym}
\DeclareMathOperator{\GL}{GL}
\DeclareMathOperator{\im}{im}
\DeclareMathOperator{\conv}{conv}
\DeclareMathOperator{\aff}{aff}
\DeclareMathOperator{\sgn}{sgn}
\newcommand{\NN}{\mathbb Z_{\ge0}}
\newcommand{\CC}{\mathbb C}
\newcommand{\SSS}{\mathfrak S}
\newcommand{\Schur}{\mathbf S}
\newcommand{\st}{\operatorname{st}}
\newcommand{\HW}{\operatorname{HW}}
\newcommand{\one}{\boldsymbol 1}
\lstdefinestyle{proofcode}{
  language=C++,basicstyle=\ttfamily\scriptsize,
  numbers=left,numberstyle=\tiny\color{gray},numbersep=6pt,
  keywordstyle=\bfseries,commentstyle=\color{black!65},
  stringstyle=\color{black},showstringspaces=false,
  breaklines=true,breakatwhitespace=false,columns=fullflexible,
  keepspaces=true,tabsize=2,frame=none,literate={λ}{{$\lambda$}}1,
  aboveskip=8pt,belowskip=8pt
}
\title{Foulkes' conjecture for the sixth symmetric power}
\author{OpenAI}
\date{September 25, 2026}
\begin{document}
\maketitle
\begin{abstract}
We prove the sixth-symmetric-power case of Foulkes' conjecture. For every
integer \(b\ge6\) and every finite-dimensional complex vector space \(V\),
there is a \(\GL(V)\)-equivariant injection
\(\Sym^6(\Sym^b V)\hookrightarrow\Sym^b(\Sym^6 V)\).
\end{abstract}
\section{The result and its context}

Foulkes' conjecture compares two ways of composing symmetric powers.
For a finite-dimensional complex vector space \(V\) and integers
\(1\le a\le b\), it asks for a \(\GL(V)\)-equivariant injection
\[
             \Sym^a(\Sym^b V)\hookrightarrow\Sym^b(\Sym^a V).
\]
Over \(\CC\), complete reducibility makes this equivalent to an
inequality between the multiplicities of every irreducible polynomial
representation. Write \(\Schur_\lambda V\) for the Schur module indexed
by a partition \(\lambda\), and \(s_\lambda\) for its character.
If \(h_r\) denotes the complete homogeneous symmetric function, then
\(h_a[h_b]\) is the stable character of \(\Sym^a(\Sym^b V)\), where
brackets denote plethysm. The conjecture says that
\(h_b[h_a]-h_a[h_b]\) is Schur-positive: all its Schur coefficients
are nonnegative integers. We prove the sixth case.

\begin{theorem}\label{thm:main}
For every integer \(b\ge6\) and every finite-dimensional complex vector
space \(V\), there is a \(\GL(V)\)-equivariant injection
\[
                 \Sym^6(\Sym^b V)\ \hookrightarrow\ \Sym^b(\Sym^6 V).
\]
Equivalently, \(h_b[h_6]-h_6[h_b]\) is Schur-positive for every \(b\ge6\).
\end{theorem}

Theorem~\ref{thm:main} resolves this case of Foulkes' conjecture
positively, with no restriction on \(\dim V\). The diagonal case
\(b=6\) is immediate; the content is the comparison for every larger
\(b\).

\paragraph{History and the canonical map.}
The conjecture originates in Foulkes' study of concomitants
\cite{Foulkes1950}. For binary forms, classical Hermite reciprocity
gives an isomorphism between the two symmetric powers for all \(a,b\);
see \cite[Introduction]{RaicuSamWeyman2022}. In higher dimension one
seeks a multiplicity inequality instead.
The case \(a=1\) is immediate. Thrall's decomposition formulas imply
\(a=2\) \cite{Thrall1942}; see also \cite[p.~237]{DentSiemons2000}.
Dent and Siemons proved \(a=3\) \cite{DentSiemons2000}.

A natural approach uses the canonical Foulkes--Howe map
\[
 \Psi_{a,b,V}:\Sym^a(\Sym^b V)\longrightarrow\Sym^b(\Sym^a V).
\]
Realize the source as symmetric tensors in \(V^{\otimes ab}\), arrange
the tensor positions in \(a\) rows of length \(b\), transpose the array,
and symmetrize the resulting blocks. McKay proved that injectivity
for every \(V\) at a seed \((a,b_0)\), with \(b_0\ge a\), propagates to every
\((a,b)\) with \(b\ge b_0\)
\cite{McKay2008,Ikenmeyer2015}.
The computation at \((4,4)\) by M\"uller and Neunh\"offer therefore
yields \(a=4\) \cite{MuellerNeunhoeffer2005};
Cheung, Ikenmeyer, and Mkrtchyan's injectivity computation at \((5,6)\)
yields \(a=5\), with the diagonal comparison again immediate
\cite[preprint, Theorem~6(a)]{CheungIkenmeyerMkrtchyan2017}.
Canonical injectivity is stronger than Foulkes' multiplicity comparison.
The canonical maps at \((5,5)\) and \((6,6)\) have nonzero kernels in
suitable dimensions, although their source and target representations
coincide \cite{MuellerNeunhoeffer2005,CheungIkenmeyerMkrtchyan2017}.
For the sixth diagonal this already occurs when \(\dim V\ge6\)
\cite[preprint, Theorem~6(c)]{CheungIkenmeyerMkrtchyan2017}.
Our finite comparison does not require injectivity of a prescribed map.

Exact character calculations give another approach.
Evseev, Paget, and Wildon verified Foulkes' multiplicity inequalities
when the two parameters sum to at most \(19\), including the sixth
case through \(b=13\)
\cite[Corollary~5.2]{EvseevPagetWildon2014}.
Their character-deflation recurrence is a predecessor of the exact
source calculation used here. The present certificates combine such
calculations with bounds on entire families of multiplicities, leaving
only a small band for exact comparison.

Brion proved eventual truth for fixed \(a\)
\cite[Corollary~1.3]{Brion1993} and later gave a dimension-dependent
bound for surjectivity of the canonical map in the opposite direction
\cite[Theorem~3.3]{Brion1997}.
That map, \(\Psi_{b,a,V}\), is realized as multiplication in the invariant
ring associated with the normalization of the Chow variety
\cite[Section~7]{Landsberg2015}; \cite[Section~2]{RaicuSamWeyman2022}.
The companion paper proves its surjectivity for \(b\ge a(a-1)\),
uniformly in dimension \cite[Theorem~1]{OpenAIQuadratic2026}.
At \(a=6\), an invariant complement to its kernel supplies an
embedding for \(b\ge30\). The remaining task is to compare all
multiplicities for \(6\le b\le29\).

For binary forms, recent generalized comparisons concern plethysms
with different parameter pairs: Raicu, Sam, Weyman, and Yang prove
maximal rank under divisibility conditions
\cite[Theorem~1.4]{RaicuSamWeymanYang2026}, and Gangl, Guti\'errez,
and Szwej identify the dual map through a substitution construction
\cite[Theorem~1.2]{GanglGutierrezSzwej2026}.
These results concern a two-dimensional underlying space; the theorem
above treats every finite dimension.

\paragraph{Proof strategy.}
The source \(\Sym^6(\Sym^b V)\) has Schur support of length at most
six. Let \(F(\lambda)\) and \(Q(\lambda)\) denote its multiplicity and
the corresponding target multiplicity, respectively, for a partition
\(\lambda\) of \(6b\). Section~\ref{sec:reduction} reduces the problem
to these six-part partitions and supplies rectangular highest-weight
shifts of known comparisons and a chart certificate
for partitions with \(\lambda_2+\cdots+\lambda_6\le b\).
The canonical multiplication map also proves the main large-degree
range in Corollary~\ref{cor:quadratic-range}.

The finite argument compares upper bounds for \(F\) with lower bounds
for \(Q\). Section~\ref{sec:upper} bounds source multiplicities by
counting possible intermediate partitions in a signed strip recurrence,
and gives exact recurrences for both characters.
Section~\ref{sec:lower} amplifies known target multiplicities:
a basis of highest-weight polynomials can be chosen with distinct leading
monomials, so a sumset bound supplies many independent products of those
polynomials.
Multiplication by a single nonzero highest-weight polynomial then
transfers these lower bounds to further weights.
Section~\ref{sec:certificates} organizes the resulting tests into a
finite computation. It checks every partition for \(6\le b\le25\).
In increasing degrees from \(26\) onward, rectangular shifts reduce
the remaining problem to \(\lambda_6<6\); the certificates leave only
partitions in the band \(\lambda_3+\cdots+\lambda_6\le27\).
Section~\ref{sec:band} computes exact differences throughout this band
by reusing seven truncated rational-function numerators.
Section~\ref{sec:verification} proves the arithmetic bounds and
assembles the range coverage.

An independent route is retained in full. The chart calculation used
in the finite comparison, combined with polarization, proves canonical
surjectivity for \(b\ge150\) in Proposition~\ref{prop:large}.
The polarization argument is related to the successive weight shifts
in the proof of McKay's theorem \cite[Section~5]{Ikenmeyer2015}.
Both certificate programs run through \(b=149\), so the extended finite
verification and this chart bound prove Theorem~\ref{thm:main} without
the quadratic-stabilization input. The main route needs only the
finite comparisons through \(b=29\).
The complete programs are printed in the appendices and supplied with
reference outputs and a reproduction script. The formulas explain
what the programs compute; their exact execution supplies the finite
nonnegativity assertions.

Throughout, we use characteristic-zero complete reducibility and
highest-weight theory for polynomial \(\GL\)-representations
\cite[Theorem~9.19 and Proposition~14.13]{FultonHarris1991},
Schur functors and the Cauchy decomposition
\cite[Theorem~6.3 and Exercise~6.11(b)]{FultonHarris1991},
and the Schur alternant formula \cite[Chapter~I, (3.1)]{Macdonald1995}.
Stable Schur coefficients are unchanged on specializing extra character
variables to zero whenever the partition length remains admissible
\cite[Chapter~I, (3.2)--(3.3)]{Macdonald1995}.
The character interpretation of plethysm is recalled in
\cite[Chapter~I, Appendix~A, Sections~7--8]{Macdonald1995}.
We also write \(e_r\) for the elementary symmetric function and set
\(h_0=e_0=1\).

\section{The large-degree range and finite reductions}\label{sec:reduction}

The proof has a representation-theoretic reduction and a finite
multiplicity comparison. We first specify the canonical map supplied by
the companion and reduce the remaining comparison to partitions with at
most six parts. A chart of the same invariant ring gives one of the finite
certificates. The final subsection proves the independent large-degree
bound used with the extended computation.

\subsection{The multiplication map}
For the moment let \(a\ge2\), and set
\[
 P=(\Sym V)^{\otimes a},\qquad
 W_j=P_{(j,\ldots,j)},\qquad
 R_j=W_j^{\SSS_a},\qquad R=\bigoplus_{j\ge0}R_j.
\]
The symmetric group permutes the tensor factors. We identify
\(\Sym^a W\) with \((W^{\otimes a})^{\SSS_a}\) by the normalized average
\[
 w_1\cdots w_a\longmapsto
 \frac1{a!}\sum_{\sigma\in\SSS_a}
 w_{\sigma(1)}\otimes\cdots\otimes w_{\sigma(a)}.
\]
Its inverse is the restriction of the symmetric-product quotient.
Let \(A\subset R\)
be the graded subalgebra generated by \(R_1=\Sym^a V\).
For \(v\in V\), put \(z(v)=v^{\otimes a}\in R_1\).
Multiplication defines the canonical Foulkes--Howe map
\begin{equation}\label{eq:multmap}
       \mu_{a,b,V}:\Sym^b(\Sym^a V)\longrightarrow
       R_b=\Sym^a(\Sym^b V),
\end{equation}
whose image is \(A_b\). Thus \(\mu_{a,b,V}=\Psi_{b,a,V}\) in the
notation of the introduction. It is equivariant, so \(A_b=R_b\) suffices
for the desired embedding, by complete
reducibility. This is the equal-multidegree invariant construction
underlying the Chow normalization; see \cite[Section~2]{RaicuSamWeyman2022}.
We will use its explicit polynomial realization.

\begin{corollary}[The main large-degree range]\label{cor:quadratic-range}
For every finite-dimensional complex \(V\) and every \(b\ge30\), there
is a \(\GL(V)\)-equivariant injection
\(\Sym^6(\Sym^b V)\hookrightarrow\Sym^b(\Sym^6 V)\).
\end{corollary}
\begin{proof}
The map in~\eqref{eq:multmap}, with exactly the averaging convention
above, is the canonical multiplication map of
\cite[Theorem~1]{OpenAIQuadratic2026}. That theorem gives \(A_b=R_b\)
for \(a=6\) and \(b\ge6(6-1)=30\), including dimension zero.
In positive dimension its kernel has an invariant complement, by
complete reducibility of complex polynomial \(\GL(V)\)-representations.
Restriction to this complement is an isomorphism onto \(R_b\);
its inverse gives an equivariant section and hence the stated injection.
There is no assertion that this section is a prescribed canonical map.
\end{proof}

\subsection{Six variables and rectangular shifts}
For the finite comparison take \(a=6\). If \(b\ge0\) and
\(\lambda\) is a partition of \(6b\), the Cauchy decomposition of
\(P=\Sym(V\otimes\CC^6)\)
\cite[Exercise~6.11(b)]{FultonHarris1991} gives the source multiplicity
\begin{equation}\label{eq:label-multiplicity}
 F(\lambda)=\dim\left((\Schur_\lambda\CC^6)_{(b,\ldots,b)}^{\SSS_6}\right).
\end{equation}
In particular it vanishes for partitions of length greater than six.
Stable Schur coefficients allow us to work in \(V=\CC^6\): extra
variables are specialized to zero, while longer target constituents
already have nonnegative multiplicities. This covers every dimension.

Put \(S=\Sym(\Sym^6\CC^6)\) with its ordinary symmetric-algebra grading,
and let \(Q(\lambda)\) be the multiplicity of \(\Schur_\lambda\CC^6\)
in \(S_b\). Pad every partition by zeros to six parts and write
\[
 d_i=\lambda_i-\lambda_{i+1},\quad \lambda_7=0,\qquad
 \lambda(d)_i=\sum_{j=i}^6d_j,\qquad
 w(d)=\sum_{i=1}^6 i\,d_i.
\]
Also put \(t_s(d)=\sum_{j>s}\lambda(d)_j\). We abbreviate
\(F(d)=F(\lambda(d))\) and \(Q(d)=Q(\lambda(d))\) when \(6\mid w(d)\);
then \(b=w(d)/6\). Gap-vector inequalities are componentwise, and \(e_i\)
denotes the \(i\)-th unit vector in gap coordinates.

\begin{lemma}\label{lem:rectangles}
Let \(c,d\in\NN^6\) satisfy \(6\mid w(c)\) and \(6\mid w(d)\).
If \(Q(c)>0\), then \(Q(d+c)\ge Q(d)\). In particular
\[
          Q(d+6e_i)\ge Q(d)\quad(1\le i\le6),
          \qquad F(d+6e_6)=F(d).
\]
\end{lemma}
\begin{proof}
Multiplication by one nonzero highest-weight polynomial of weight
\(\lambda(c)\) injects the highest-weight space at \(\lambda(d)\)
into that at their sum, because \(S\) is a domain.
This is the usual highest-weight multiplication argument; compare
\cite[preprint, Lemma~2.2]{BurgisserIkenmeyerPanova2019}.

The rectangular occurrence \(Q(6e_i)>0\) is a special case of the
even-partition theorem of B\"urgisser, Christandl, and Ikenmeyer
\cite[Section~1.1, Theorem]{BurgisserChristandlIkenmeyer2011}.
We give the determinant-tensor construction of
\cite[preprint, Corollary~4.8]{BurgisserIkenmeyerPanova2019}
in the form needed here. Take six copies of the determinant tensor on
the first \(i\) basis vectors. Regroup the tensor positions into \(i\)
blocks, each taking the same position from every copy, and symmetrize
to \(\Sym^i(\Sym^6\CC^6)\). This is a highest vector of weight
\((6,\ldots,6)\) with \(i\) parts. Pair each block with the same sum of
pure sixth powers of the first \(i\) dual basis vectors. A surviving
term uses the same permutation in all six determinant copies, so its
sign is \(+1\); such terms exist. The vector is nonzero.

Adding \(6e_6\) adds six to every partition part.
In~\eqref{eq:label-multiplicity}, this tensors with \(\det^6\),
shifts the equal label weight by six, and is trivial on permutation
matrices. Thus \(F(d+6e_6)=F(d)\).
\end{proof}

It follows that, after checking all partitions for \(6\le b\le25\),
induction in \(26\le b\le29\) need only treat \(d_6<6\).
If \(d_6\ge6\), subtract \(6e_6\); the degree becomes \(b-6\ge20\),
and Lemma~\ref{lem:rectangles} imports the earlier comparison.
The same induction works throughout \(26\le b\le149\), the full
range retained in the programs for the independent route.

\subsection{A chart certificate for the first row}
We return to arbitrary \(a\ge2\) and \(V\) for the chart calculation.
Its factor bound will show, at \(a=6\), that the multiplication image
contains every source highest-weight space with \(\lambda_1\ge5b\).
It will also provide the independent large-degree bound in the next
subsection. Dimension zero is immediate, so assume \(V\ne0\).

Fix \(0\ne v\in V\), use it as the first basis vector, and write
\(x_{i0},x_{i1},\ldots\) for the corresponding polynomial variables in
row \(i\) of \(P\). Put \(z=z(v)=\prod_i x_{i0}\).
After inverting \(z\), every \(x_{i0}\) is invertible. Write
\(Y_i=(x_{i1}/x_{i0},x_{i2}/x_{i0},\ldots)\).

\begin{lemma}\label{lem:chart}
For every integer \(t\ge0\), every row-permutation-invariant polynomial
in the lists \(Y_i\) of total degree at most \(t\) is a linear combination of products of at
most \(t\) elements of \(z^{-1}R_1\).
Consequently, for integers \(k,L\ge0\) and \(f\in R_k\),
\[
             z^Lf\in A_{k+L}\qquad\text{if }k+L\ge ak.
\]
\end{lemma}
\begin{proof}
The multisymmetric generators are described in
\cite[Remark~2.9]{RaicuSamWeyman2022}; we include the factor-length
argument needed here.
For the general multisymmetric generator theorems, see
\cite[Theorem~1]{Vaccarino2005} and \cite[Corollary~8.4]{Rydh2007}.
For each linear form \(\ell\) in one list, the coefficients of
\[
 \frac{z(v+tu)}{z}=\prod_{i=1}^a(1+t\ell(Y_i))
\]
show that \(e_h(\ell(Y_1),\ldots,\ell(Y_a))\), and all their
polarizations, lie in \(z^{-1}R_1\), for \(1\le h\le a\).
Here \(u\) ranges over the span of the remaining basis vectors.

Give \(e_h\) weight \(h\). Newton identities express
\(\sum_i\ell(Y_i)^d\) as a polynomial in \(e_1,\ldots,e_a\)
of weighted degree \(d\), including when \(d>a\). Each product has
at most \(d\) factors. Polarizing in \(\ell\) therefore expresses
\[
                         p_m=\sum_i m(Y_i)
\]
for every degree-\(d\) monomial \(m\) using at most \(d\) polarized
elementary factors.

Invariant monomial orbit sums are scalar multiples of expressions
\[
        \sum_{\substack{i_1,\ldots,i_h\\\text{distinct}}}
                    \prod_{\nu=1}^h m_\nu(Y_{i_\nu}),
\]
where the \(m_\nu\) have positive degree. Inclusion-exclusion over
collisions of the indices expresses this in products of \(p_m\)'s;
a collision replaces monomials by their product and preserves total
degree. Thus total degree at most \(t\) requires at most \(t\) elementary
factors. Constants require no factors.

For \(f\in R_k\), the polynomial \(f/z^k\) is invariant and has total
list degree at most \(ak\). A term with \(m\le ak\) ratio factors has
the form \(z^{-m}g_1\cdots g_m\), with \(g_\nu\in R_1\).
Multiplication by \(z^{k+L}\) gives an element of \(A_{k+L}\) when
\(k+L\ge ak\). Equality in the localization descends to \(P\), a domain.
\end{proof}

\begin{corollary}[First-row certificate]\label{cor:first-row}
For every integer \(b\ge1\) and partition \(\lambda\) of \(6b\) with
\(\sum_{i\ge2}\lambda_i\le b\), one has
\[
             [s_\lambda]h_6[h_b]\le [s_\lambda]h_b[h_6].
\]
\end{corollary}
\begin{proof}
Partitions of length greater than six have zero source multiplicity, so
work in \(V=\CC^6\) with its standard ordered basis and put \(a=6\).
Choose \(v\) to be the first basis vector in the chart above.
For every weight-\(\lambda\) polynomial \(f\in R_b\), the total
list degree of \(f/z^b\) is
\(t=\lambda_2+\cdots+\lambda_6\).
The first assertion of Lemma~\ref{lem:chart} writes this ratio as a
linear combination of terms \(z^{-m}g_1\cdots g_m\), with
\(m\le t\le b\) and \(g_\nu\in R_1\). Thus every corresponding
term of \(f\) is
\[
                     z^{b-m}g_1\cdots g_m\in A_b.
\]
The entire weight-\(\lambda\) space of \(R_b\) therefore lies in
\(A_b\). Since \(A_b\) is a \(\GL(V)\)-submodule of \(R_b\), their
highest-weight spaces at \(\lambda\) are equal: intersect the common
weight space with the kernels of the positive simple-root operators.
The surjection onto \(A_b\) induced by~\eqref{eq:multmap} then gives
\(Q(\lambda)\ge F(\lambda)\).
\end{proof}

\subsection{An independent large-degree bound}
We now prove surjectivity of \(\mu_{a,b,V}\) for
\(b\ge a(a-1)^2\), without using the companion. The chart will clear
denominators uniformly once we know that \(R\), as an \(A\)-module,
is generated in degrees below \(a\). Polarization supplies that
finite-generation statement and then turns the chart calculation into
surjectivity in large degree. We continue with arbitrary \(a\ge2\);
the zero-dimensional case is immediate.

\begin{lemma}[Polarization]\label{lem:polarization}
For integers \(L\ge1\) and \(j\ge aL\),
\[
       W_j=\sum_{v\in V}z(v)^L W_{j-L},
       \qquad
       R_j=\sum_{v\in V}z(v)^L R_{j-L}.
\]
The sums mean linear spans.
\end{lemma}
\begin{proof}
The operator argument is related to the successive weight shifts in
\cite[Section~5]{Ikenmeyer2015}. Polarize the degree-\(aL\) polynomial
\(v\mapsto z(v)^L\), and multiply, to obtain a linear map
\[
                 \Sym^{aL}V\otimes W_{j-L}\longrightarrow W_j.
\]
Its image is the first displayed span. Identify \(W_j^*\) with
polynomials \(f(x_1,\ldots,x_a)\), separately homogeneous of degree \(j\).
Introduce one auxiliary vector variable \(y\). The dual map is, up to a nonzero scalar,
\[
                f\longmapsto\prod_{i=1}^a
                       (y\cdot\partial_{x_i})^L f.
\]
Indeed, evaluation of the dual pairing extracts the coefficient of
\(t_1^L\cdots t_a^L\) in \(f(x_1+t_1y,\ldots,x_a+t_ay)\).

On the polynomial space of total degree \(aj\) in
\(x_1,\ldots,x_a,y\), the operators
\[
 D_i=y\cdot\partial_{x_i},\quad
 E_i=x_i\cdot\partial_y,\quad
 H_i=\deg(x_i)-\deg(y)
\]
form an \(\mathfrak{sl}_2\), with \(D_i\) lowering weight by two.
The other lists are spectator variables. The standard weight strings
are described in \cite[Section~11.1, (11.5)]{FultonHarris1991}.
In an irreducible module of highest weight \(n\), a weight \(w\ge L\)
is at distance \((n+w)/2\ge L\) from the bottom. Therefore \(D_i^L\)
is injective on the full weight-\(w\) space of any finite-dimensional
module. After the first \(i-1\) shifts, the relevant weight is
\(j-(i-1)L\ge L\). Every shift is injective on a space containing the
preceding image, so their product is injective. Duality proves the
first identity. Average its coefficients over \(\SSS_a\); each
\(z(v)^L\) is invariant, proving the second identity.
\end{proof}

\begin{corollary}\label{cor:module-generation}
The \(A\)-module \(R\) is generated by \(R_0,\ldots,R_{a-1}\).
\end{corollary}
\begin{proof}
Lemma~\ref{lem:polarization} with \(L=1\) gives
\(R_j=R_1R_{j-1}\) for \(j\ge a\); apply induction on \(j\).
\end{proof}

\begin{proposition}[Independent chart bound]\label{prop:large}
For \(a\ge2\), one has \(R_b=A_b\) whenever
\(b\ge a(a-1)^2\).
\end{proposition}
\begin{proof}
Take \(L=(a-1)^2\), so that \(L\ge(a-1)k\) for every
\(0\le k<a\). Lemma~\ref{lem:chart} therefore applies to each of
these degrees. For any fixed nonzero \(v\), multiply an \(A\)-linear
expression in the generators of Corollary~\ref{cor:module-generation}
by \(z(v)^L\); this gives \(z(v)^LR\subset A\).
The same assertion is trivial for \(v=0\). If \(b\ge aL\),
Lemma~\ref{lem:polarization} now yields
\[
                  R_b=\sum_v z(v)^L R_{b-L}\subseteq A_b.
\]
The opposite inclusion is definitional.
\end{proof}

For \(a=6\), Proposition~\ref{prop:large} supplies the independent
large-degree range \(b\ge150\). Combined with the retained finite
verification through \(b=149\), it proves Theorem~\ref{thm:main}
without the quadratic-stabilization input. The main proof needs only
\(6\le b\le29\), since Corollary~\ref{cor:quadratic-range} applies
from \(b=30\).

\section{An upper bound and two exact recurrences}\label{sec:upper}

We compare the source and target multiplicities without computing both
at every partition. This section supplies a gap-monotone upper bound
\(U(d)\ge F(d)\), together with exact recurrences for the target tables
and for the source values needed when that bound is inconclusive.
All three constructions come from coefficient extraction with Schur
alternants.

\subsection{The signed strip rule}
The coefficient rule below is the complete-symmetric case of the
plethystic Murnaghan--Nakayama rule; see
\cite[pp.~28--29]{DesarmenienLeclercThibon1994} and
\cite[arXiv:1408.3554v2, equation~(2)]{Wildon2016}.
Its residue-class determinant formulation also appears in
\cite[Theorem~3.4 and Corollary~3.11]{CaoJingLiu2025}.
We include the proof and the containment inequalities in the form
needed by the exact calculation.

For \(n\) variables put \(\delta_n=(n-1,\ldots,0)\).
The next rule bounds a single removal step in a product of factors
\(h_b(X^i)\), and also enumerates that step exactly.

\begin{lemma}[Signed strips]\label{lem:strips}
Let \(\gamma,\mu\) be partitions of length at most \(n\), and suppose
\(|\gamma|-|\mu|=ib\), where \(i,b\ge1\). Then
\[
                  [s_\gamma]\,h_b(X^i)s_\mu(X)\in\{0,1,-1\}.
\]
For a nonzero coefficient one necessarily has
\[
        \gamma_j\ge\mu_j,\qquad \mu_j\ge\gamma_{j+i}
                         \quad\text{whenever the subscripts are in range}.
\]
Here \(X^i=(x_1^i,\ldots,x_n^i)\).
\end{lemma}
\begin{proof}
Set \(l=\gamma+\delta_n\) and \(a=\mu+\delta_n\).
Multiplying by the Schur alternant denominator
\cite[Chapter~I, equation~(3.1)]{Macdonald1995} gives
\[
 [s_\gamma]h_b(X^i)s_\mu(X)
 =\det\bigl[\one_{\{a_k\le l_j,\ a_k\equiv l_j\pmod i\}}\bigr]_{j,k=1}^n.
\]
Indeed its determinant expansion assigns the entries of \(a\) to
positions \(j\), with assigned values \(\beta_j\) satisfying
\[
                  \beta_j\le l_j,\qquad \beta_j\equiv l_j\pmod i.
\]
The size condition makes the sum of the decrement quotients equal to
\(b\), so there is no further constraint from \(h_b(X^i)\).

Group the rows and columns by residue modulo \(i\), preserving decreasing
order within each group. Unequal row and column counts in a residue give
zero, and empty blocks contribute one. In a square block of size \(m\),
each row support is a terminal interval of columns, and the supports
decrease with the row index. A zero row or repeated support makes the
determinant zero. Otherwise their sizes are \(m,m-1,\ldots,1\): the block
is upper triangular with diagonal entries one, and its only allowed
matching is the diagonal. Translating this matching back to the original
positions gives
\begin{equation}\label{eq:stripconditions}
 \beta_j\le l_j,\qquad \beta_j>l_{j'}
      \quad\text{if \(j'\) is the next position of the same residue}.
\end{equation}
There is therefore just one surviving assignment. The sign from sorting
the assembled \(\beta\)'s is the remaining row-and-column ordering sign,
which proves the asserted coefficient values.

Sorting the componentwise inequalities \(\beta_j\le l_j\) gives
\(\mu_j\le\gamma_j\). For the other inequality, consider the first \(j+i\)
entries of \(l\). At most one per residue lacks a preceding entry of
that residue. Each of the remaining at least \(j\) entries has a
predecessor position whose assigned value, by
\eqref{eq:stripconditions}, is at least \(l_{j+i}+i\).
The \(j\)-th largest assigned value is therefore at least
\(l_{j+i}+i\), which after removing the staircases says
\(\mu_j\ge\gamma_{j+i}\).
\end{proof}

\subsection{A monotone upper certificate}
Averaging the permutations of six tensor factors projects onto
\(\Sym^6(\Sym^b V)\). Taking traces gives the cycle formula
\begin{equation}\label{eq:cycle}
        h_6[h_b](X)=\frac1{720}\sum_{\pi\in\SSS_6}
                   \prod_{\text{cycles }c\text{ of }\pi}h_b(X^{|c|}).
\end{equation}
We bound a coefficient of each product by removing its factors one at
a time. The signed strip rule restricts the intermediate partitions;
the bound below counts their possible coordinates while discarding
their signs and some of their constraints.

For a nonnegative gap vector \(d\), define
\[
 h_{sj}=1+\lambda(d)_j-\lambda(d)_{j+6-s},\qquad
 N_s(d)=\frac{\prod_{j=1}^s h_{sj}}{\max_{1\le j\le s}h_{sj}}
            \quad(2\le s\le5),\qquad N_0=N_1=1.
\]
For an ordered list \(I=(i_1,\ldots,i_v)\) with positive entries summing
to six put
\[
 B_I(d)=\prod_{u=1}^v N_{6-i_1-\cdots-i_u}(d),\qquad
 U(d)=\left\lfloor\frac1{720}
       \sum_{\pi\in\SSS_6}\min_{I\sim\pi}B_I(d)\right\rfloor ,
\]
where \(I\sim\pi\) means any ordering of the cycle lengths of \(\pi\).
This definition makes sense even when \(6\nmid w(d)\).

\begin{proposition}\label{prop:upper}
The function \(U\) is nondecreasing in every gap. For every nonnegative
gap vector \(d\) with \(6\mid w(d)\), one has \(F(d)\le U(d)\).
\end{proposition}
\begin{proof}
Each \(N_s\) is the minimum of the products obtained by omitting one of
the widths \(h_{sj}\). Every width is nondecreasing in the gaps, proving
monotonicity through the products, minima, sum, and floor.

The zero vector has \(F(0)=U(0)=1\), so assume \(b=w(d)/6\ge1\).
Remove the factors of one product in \eqref{eq:cycle} in an order \(I\).
After lengths totaling \(6-s\) have been removed, the remaining product
is the trace of a permutation of the \(s\) factors on
\((\Sym^b V)^{\otimes s}\), together with the diagonal \(\GL(V)\)-action.
The permutation commutes with that action, so its Schur support lies
among the constituents of the tensor power. Those have length at most
\(s\) by Pieri's rule
\cite[Chapter~I, equation~(5.16)]{Macdonald1995}.
Iterating Lemma~\ref{lem:strips}, every possible intermediate
partition \(\mu\) satisfies
\[
       |\mu|=sb,\qquad
       \lambda_{j+6-s}\le\mu_j\le\lambda_j\quad(1\le j\le s).
\]
Fix all but one coordinate; the sum fixes the remaining one.
Thus there are at most \(N_s(d)\) choices for \(2\le s\le5\), and at most
one for \(s=0,1\). Each chain has absolute coefficient at most one.
The coefficient of the product consequently has absolute value at most
\(B_I(d)\), for every order \(I\). Average these bounds in
\eqref{eq:cycle} and use the integrality of \(F(d)\).
\end{proof}

In Appendix~\ref{app:certificate}, \texttt{init\_masks} lists the eleven
cycle types with their class sizes and the remaining lengths
\(s\in\{2,3,4,5\}\) for each removal order. The routine
\texttt{ub\_cycle} forms the corresponding products of \(N_s\), takes
their minima, and performs the weighted sum and division by \(720\).

\subsection{Exact coefficients}
The target tables and the remaining source comparisons use Newton's
identity in place of a bound. For an integer vector \(v\) of length
\(n\), a \emph{signed lookup} is zero if \(v\) has a negative entry or a
repeated entry. Otherwise sort \(v\) decreasingly, subtract
\(\delta_n\), and multiply the coefficient at that partition by the
sorting sign. A strictly decreasing nonnegative integer list is at
least \(\delta_n\) componentwise, so this subtraction always gives a
partition.

Write \(Q_b(\lambda)=[s_\lambda]h_b[h_6]\) and
\(\widetilde Q_b(v)\) for its signed lookup in \(n\) variables.
The initial value is
\(Q_0(0)=1\). Newton's identity
\cite[Chapter~I, equation~(2.11)]{Macdonald1995}, multiplied by the alternating
denominator, gives
\begin{equation}\label{eq:targetrec}
 bQ_b(\lambda)=
 \sum_{i=1}^b\ \sum_{\substack{\alpha\in\NN^n\\|\alpha|=6}}
       \widetilde Q_{b-i}(\lambda+\delta_n-i\alpha).
\end{equation}
Indeed \(h_6(X^i)=\sum_{|\alpha|=6}X^{i\alpha}\).
Every nonzero previous partition \(\mu\) in this formula satisfies
\(\mu\le\lambda\) componentwise: subtracting a nonnegative vector cannot
increase any decreasing order statistic, and the same staircase is then
subtracted. In particular, a bound on a partition's tail is preserved in
all previous lookups.

For the source, fix \(b\) and set
\[
 H_j(\gamma)=[s_\gamma]h_j[h_b],\qquad H_0(0)=1.
\]
Newton's identity now gives the source recurrence; compare the
character formulation in
\cite[Proposition~5.1]{EvseevPagetWildon2014}:
\begin{equation}\label{eq:sourcerec}
 jH_j(\gamma)=\sum_{i=1}^j\ \sum_{\mu}
       [s_\gamma]\bigl(h_b(X^i)s_\mu(X)\bigr)H_{j-i}(\mu),
       \qquad |\mu|=(j-i)b.
\end{equation}
At stage \(j\), \(n=j\) variables suffice. Lemma~\ref{lem:strips} gives
an economical enumeration. Put \(l=\gamma+\delta_j\).
In each residue modulo \(i\), its last assigned
value must be that residue's least nonnegative representative: the
partition \(\mu\) has length at most \(j-i\), so its shifted list ends in
\(i-1,\ldots,0\). If any residue is absent, there is no term. At every
other position \(h\), whose next position in its residue is \(h'\), choose
\[
 \beta_h=l_h-ix_h,\qquad
               0\le x_h\le(l_h-l_{h'})/i-1.
\]
The sum of all decrement quotients, including the fixed last ones, is
\(b\). These choices have distinct values: different residues cannot
coincide, and in the same residue
\(\beta_h>l_{h'}\ge\beta_{h'}\). The fixed representatives
\(0,\ldots,i-1\) are the smallest \(i\) values. Thus sorting and
subtracting \(\delta_j\) gives a partition \(\mu\) of length at most
\(j-i\), with the sorting sign. Every such choice contributes once.
This is exactly the routine \texttt{strip} in Appendix~\ref{app:certificate}.
Equation~\eqref{eq:targetrec} is implemented by \texttt{gen},
\texttt{Make}, and \texttt{QTail}; \eqref{eq:sourcerec} by \texttt{Source}.

\section{A power bound from leading monomials}\label{sec:lower}

Rectangular shifts give a first lower bound for \(Q\). The following
bound amplifies a known multiplicity.

The sumset estimate is the repeated-set case of the Matolcsi--Ruzsa
inequality \cite[preprint, Theorem~1.5]{MatolcsiRuzsa2010};
take \(A=B=D\) and \(r=k-1\) there when \(k\ge2\).
See also \cite[author version, Corollary~2]{BoroczkySantosSerra2014}.
We include a self-contained triangulation proof. Its generic-point
assignment is closely related to the half-open-simplex count used in
\cite[author version, Theorem~20]{BoroczkySantosSerra2014}, where the
regions are instead obtained from a shelling.

\begin{lemma}[Sumsets]\label{lem:sumset}
Let \(D\) be a nonempty finite set of \(C\) points in a real vector space,
of affine dimension \(h\). Write
\(kD=\{v_1+\cdots+v_k:v_1,\ldots,v_k\in D\}\).
For this \(k\)-fold sumset, \(k\ge1\), one has
\[
 |kD|\ \ge\
 f(C,h,k):=\binom{k+h}{h}+(C-1-h)\binom{k+h-1}{h}.
\]
\end{lemma}
\begin{proof}
The case \(h=0\) is immediate. Triangulate \(\conv D\) using every point
of \(D\) as a vertex and no other vertices. Such a triangulation can be
built by starting with a full-dimensional simplex from \(D\), then
inserting points. For a point in the current hull, subdivide the simplices
through its minimal containing face in the current triangulation by
coning from it, retaining the induced common face subdivisions.
For a point outside the hull, cone from it to the
boundary simplices on strictly visible facets. Rays from the new point
give the intersection and coverage assertions in the latter construction.
Both operations preserve all previous vertices.

Choose \(p\) in the relative interior of the hull, off all hyperplanes
of facets of full-dimensional simplices. Assign \(x\in\conv D\) to the
unique full simplex whose interior contains \((1-\epsilon)x+\epsilon p\)
for all sufficiently small positive \(\epsilon\). There are finitely many
facet hyperplanes, so this rule is well defined.

In barycentric coordinates of one full simplex, this assigned region
requires nonnegative coordinates, with strict positivity exactly at the
indices where the corresponding coordinate of \(p\) is negative.
Assign a sum of \(k\) vertices according to its average, which lies in
\(\conv D\). If there are \(e\) strict indices, the sums of \(k\) vertices
of this simplex whose averages are assigned to it number
\(\binom{k-e+h}{h}\), interpreted as zero if \(k<e\).
Affine independence makes these points distinct; the assignment makes
the counted sets for different simplices disjoint subsets of \(kD\).

There is exactly one simplex with \(e=0\), namely the one containing
\(p\). At \(k=1\) the counting rule counts precisely the \(C\) vertices.
To see this, a vertex of the complex lying in a full simplex must be a
vertex of that simplex, by the common-face property. Its assigned
simplex therefore counts it. A simplex with \(e=0\) contributes \(h+1\),
one with \(e=1\) contributes \(1\), and all others contribute zero.
Thus exactly \(C-h-1\) simplices have \(e=1\). Keeping just the \(e=0,1\)
contributions for arbitrary \(k\) proves the bound.
\end{proof}

The leading-monomial argument in the next proposition is an elementary
instance of the value-semigroup method: dimensions are counted by
distinct values, and multiplication adds values; see
\cite[preprint, Propositions~2.3, 2.6, and~2.10]{KavehKhovanskii2012}.
We give the finite-degree argument and the dimension estimate explicitly,
since their quantitative form is used in the certificate.

\begin{proposition}[Power certificate]\label{prop:power}
Suppose \(Q(x)\ge C>0\), where \(C\) is an integer and
\(w(x)=6d_0\) with \(d_0>0\). Put
\[
 z=\min\left(18,\ \min\left\{h\ge0:
                              \binom{d_0+h}{h}\ge C\right\}\right).
\]
Then for all \(k\ge1\),
\begin{equation}\label{eq:power}
                Q(kx)\ge f(C,z,k).
\end{equation}
\end{proposition}
The cap at \(18\) only weakens the lower bound; it limits the integer
arithmetic in Section~\ref{sec:verification}.

\begin{proof}
Choose \(C\) independent highest-weight polynomials at \(x\) in the
polynomial algebra \(S=\Sym(\Sym^6\CC^6)\). Row reduction with respect to
a fixed monomial order gives a basis with distinct leading monomials.
Their exponent vectors form a set \(D\) of size \(C\), all with
nonnegative integer entries summing to \(d_0\).
For each point of \(kD\), choose a product of \(k\) basis elements with
that leading exponent. The products have distinct leading monomials
and hence are independent highest-weight polynomials at \(kx\).
Thus \(Q(kx)\ge|kD|\).

Let \(h=\dim\aff D\). Some projection onto \(h\) coordinate positions
is injective on \(\aff D\), by a nonzero maximal minor of its direction
space. The projected integer points are nonnegative with sum at most
\(d_0\), so
\[
                     C\le\binom{d_0+h}{h},\qquad z\le h\le C-1.
\]
Lemma~\ref{lem:sumset} applies. Finally, for \(0\le j<h\) one computes
\[
 \frac{f(C,j+1,k)-f(C,j,k)}{\binom{k+j-1}{j}}
                  =\frac{(k-1)(C-j-1)}{j+1}\ge0.
\]
This proves \(f(C,h,k)\ge f(C,z,k)\). The formula also covers \(C=1\),
when \(h=z=0\), and \(k=1\).
\end{proof}

\section{Certificates in degrees below 150}\label{sec:certificates}

We now combine the upper bound \(F\le U\), exact target multiplicities,
and highest-weight multiplication. The base interval \(6\le b\le25\)
is checked directly. Beyond it, exact target values supply lower-bound
seeds: they are computed from the target recurrence without assuming
the desired comparison. The resulting certificates leave a set of
partitions whose exact differences will be checked in
Section~\ref{sec:band}. We retain the full interval through \(b=149\)
for the independent chart proof; only \(b\le29\) is needed with
quadratic stabilization.

\subsection{Initial tables and the base interval}
Compute \(Q\) using \eqref{eq:targetrec} on the following four sets.
All partitions have length at most six.
\begin{equation}\label{eq:tables}
\begin{array}{c|l|c}
\text{degrees}&\text{additional restriction}&\text{variables used}\\ \hline
0\le b\le25&\text{none}&6\\
0\le b\le149&\operatorname{length}(\lambda)\le3&3\\
0\le b\le55&\operatorname{length}(\lambda)\le4&4\\
0\le b\le45&t_3(d)\le18&6
\end{array}
\end{equation}
Previous-degree lookups remain in the applicable set by
Section~\ref{sec:upper}. The first table supplies the direct base
comparisons; all four supply lower-bound seeds for larger degrees.
Overlaps are immaterial.

For \(6\le b\le25\), every partition is first tested against \(Q(d)\ge
U(d)\). When this fails, compute \(F(d)=H_6(\lambda(d))\) exactly by
\eqref{eq:sourcerec}. Only these requested values at the last stage
\(j=6\) are needed. Table~\ref{tab:base} records all the fallback
comparisons and their equalities.

\begin{table}[ht]
\centering
\begin{tabular}{r|r|r@{\qquad}r|r|r}
\(b\)&exact comparisons&equalities&\(b\)&exact comparisons&equalities\\ \hline
6&2428&2428&16&7370&354\\
7&4476&315&17&7599&369\\
8&6820&255&18&7801&384\\
9&6899&250&19&8057&399\\
10&6886&264&20&8233&414\\
11&6814&279&21&8427&429\\
12&6850&294&22&8650&444\\
13&6872&309&23&8850&459\\
14&6973&324&24&9031&474\\
15&7131&339&25&9268&489
\end{tabular}
\caption{The base interval: \(145435\) exact fallback comparisons,
\(9272\) equalities, and no negative difference. All other partitions
pass \(Q\ge U\).}\label{tab:base}
\end{table}

\subsection{Residues and transport of lower bounds}
For \(26\le b\le149\) we need only consider \(d_6<6\), by
Lemma~\ref{lem:rectangles}. There is a unique expression
\[
 d=r+6m,\qquad
 r\in\mathcal R=\{0,\ldots,5\}^6\cap\{6\mid w(r)\},\qquad
 m\in\mathcal G=\{m\in\NN^6:m_6=0,\ w(m)\le149\}.
\]
Write \(b_0(r)=w(r)/6\), so that \(b=b_0(r)+w(m)\).
There are \(7776\) residues: choose \(r_2,\ldots,r_6\), and \(r_1\) is
uniquely determined modulo six. Moreover \(b_0(r)\le17\).
The grid \(\mathcal G\) has \(6611697\) points. We call its indices
\(m\) modes: fixing \(m\) groups the \(7776\) possible residues for
simultaneous testing. Pairs whose degree lies outside
\(26\le b\le149\) will not require a test.

For an array on a downward-closed set, \emph{propagation} replaces its
value at \(m\) by the maximum of all values at indices \(u\le m\).
Extending an array means filling new positions by zero and propagating.
An increasing-\(w\) sweep taking maxima with the five immediate
predecessors performs this operation. A reverse sweep through the same
predecessors forms the downward closure of a set.
Lemma~\ref{lem:rectangles} shows that propagation preserves every lower
bound for \(Q(r+6m)\), with \(r\) fixed.

The same multiplication lemma also transports a multiplicity from
residue zero to residue \(r\). Suppose \(J(u)\le Q(6u)\) for
\(u\in\mathcal G\), \(M(m)\le Q(r+6m)\), and \(E\subseteq\mathcal G\)
is a set of indices satisfying
\(Q(r+6c)>0\) for \(c\in E\). Multiplication by one nonzero
highest-weight polynomial at \(r+6c\) gives
\begin{equation}\label{eq:transport}
 Q(r+6m)\ge
 \max\left(\{M(m)\}\cup
      \{J(m-c):c\in E,\ c\le m\}\right).
\end{equation}
The right side certifies the comparison whenever it is at least
\(U(r+6m)\). No product of two multiplicity dimensions is used: a
single nonzero factor gives each injection.

We first construct \(J\) and choices of \(M,E\) valid for every residue,
so that one test can discard a whole mode. For the modes left over we
use more precise choices depending on \(r\). The following construction
specifies all arrays and tests; the common cap \(Z=2^{100}\) is applied
only to lower bounds.

\subsection{Certificates shared by all residues}

\paragraph{A shared power table.}
Initialize \(J(m)=1\) on \(\mathcal G\), using the rectangle occurrences
and the nonzero constant at \(m=0\). For every gap vector \(x\) with
\(x_6=0\), \(Q(x)>1\), and known value in one of the first three sets
of \eqref{eq:tables}, put \(d_0=w(x)/6\) and \(C=\min(Q(x),Z)\).
For each integer \(k\ge1\) with \(kd_0\le149\) and all entries of \(kx\)
divisible by six, update
\[
             J(kx/6)\ \gets\
                 \max\{J(kx/6),\ \min(Z,f(C,z,k))\},
\]
where \(z\) is as in Proposition~\ref{prop:power}. Then propagate.
It follows that
\begin{equation}\label{eq:Jbound}
                         J(m)\le Q(6m)\quad(m\in\mathcal G).
\end{equation}

\paragraph{Common multiplicities and occurrences.}
For each \(r\), initialize the exact array \(Q(r+6u)\) on
\(w(u)\le25-b_0(r)\), then extend it to \(w(m)\le42\).
Take the pointwise minimum of these propagated arrays over all \(r\),
with a cap at \(Z\), to obtain \(M_0\) on that grid.
Let \(E_0\) be the minimal points where \(M_0>0\), before extending
\(M_0\) to all of \(\mathcal G\). Then
\begin{equation}\label{eq:Mbound}
 M_0(m)\le Q(r+6m)\quad\text{for every }r,\qquad
 Q(r+6c)>0\quad\text{for every }r\text{ and }c\in E_0.
\end{equation}
The exact entries supplying the minimum may come from different
predecessors for different \(r\). Their common consequence is precisely
the two universal inequalities in \eqref{eq:Mbound}.

\paragraph{Selecting modes.}
Keep a mode \(m\in\mathcal G\) unless either \(w(m)+17\le25\), or
\begin{equation}\label{eq:sharedtest}
 \max\left(\{M_0(m)\}\cup
       \{J(m-c):c\in E_0,\ c\le m\}\right)
                   \ge U(6m+(5,5,5,5,5,5)).
\end{equation}
The first condition puts every residue in the already settled base
range whenever its degree is at least six. For the second, use
\eqref{eq:transport} with \(M=M_0\) and \(E=E_0\).
Proposition~\ref{prop:upper} bounds \(F(r+6m)\) by
\(U(6m+(5,5,5,5,5,5))\), because \(r\le(5,5,5,5,5,5)\).
The high-residue gap vector need not be admissible; \(U\) is defined and
monotone there nonetheless.

\subsection{Tests for each remaining residue}
Form the downward closure of the retained modes. On this closure,
initialize \(Q(r+6m)\) wherever it is available in
\eqref{eq:tables}, use zero elsewhere, and propagate to obtain \(M_r\).
Let \(E_r\) be the minimal positive locations of the exact array
\(Q(r+6u)\) on \(w(u)\le25-b_0(r)\).
For each originally retained mode with \(26\le b_0(r)+w(m)\le149\),
put \(d=r+6m\) and clear the pair if either
\[
 t_1(d)\le b
 \quad\text{or}\quad
 \max\left(\{M_r(m)\}\cup
          \{J(m-c):c\in E_r,\ c\le m\}\right)\ge U(d).
\]
Flag it otherwise. The first test is Corollary~\ref{cor:first-row}.
The second is \eqref{eq:transport} with \(M=M_r\) and \(E=E_r\).

Minimal positive locations can be found by checking the immediate
predecessors. Indeed positive exact multiplicities form an upper set
within the small grid, by the rectangular shifts; the propagated common
array has the same property. The set used for \(M_r\) retains every
predecessor needed for propagation. Each difference \(m-c\) in a mixed
test is nonnegative and lies in \(\mathcal G\).

In Appendix~\ref{app:certificate}, \texttt{powerAll} stores \(J\),
\texttt{uniBase} and \texttt{uOcc} store \(M_0\) and \(E_0\).
The routine \texttt{setup} retains the modes and their downward closure,
\texttt{init\_occ} constructs \(E_r\), and \texttt{TestSpec} propagates
\(M_r\) and tests the remaining pairs. The routine \texttt{run} performs
the base-interval comparisons.

\begin{verification}\label{ver:certificates}
The complete program in Appendix~\ref{app:certificate} verifies the
base comparisons in Table~\ref{tab:base}. Its mode list has \(87354\)
points, or \(114314\) after taking the downward closure. It flags
\(4969698\) pairs in the middle interval, all with \(t_2(d)\le27\).
Their counts by degree range appear in Table~\ref{tab:band}.
\end{verification}

Every pair in \(26\le b\le149\) cleared by these tests satisfies
\(Q\ge F\), by the proved upper and lower bounds or the chart.
Verification~\ref{ver:certificates}
supplies the separate finite assertion that every pair left over has
\(t_2\le27\). Thus it remains only to check the entire band
\(t_2\le27\) in \(26\le b\le149\).

\section{Exact differences in the small-tail band}\label{sec:band}

We compute the entire band
\[
\begin{gathered}
 26\le b\le149,\qquad
 \lambda=(6b-B-|\eta|,B,\eta),\\
 \eta_1\ge\cdots\ge\eta_4\ge0,\quad
 |\eta|\le27,\quad
 \eta_1\le B\le\left\lfloor\frac{6b-|\eta|}{2}\right\rfloor .
\end{gathered}
\]
In particular it contains every flagged partition of
Verification~\ref{ver:certificates}.
For each fixed tail \(\eta\), we will construct seven Laurent series
independent of \(b\). Their coefficients at
\(B,B-b,\ldots,B-6b\) will sum to \(720(Q(\lambda)-F(\lambda))\).
The same seven rows therefore serve every degree and every second part
in the displayed range.
Write \(\mathcal Q_b=h_b[h_6]\), \(\mathcal F_b=h_6[h_b]\), and set
\[
 \mathcal C(q,y)=(1-q)\prod_{\ell=1}^4(1-y_\ell)(1-y_\ell/q),\quad
 \mathcal A(y)=\prod_{\ell<h}(1-y_h/y_\ell),\quad
 D(q)=\prod_{h=1}^6(1-q^h).
\]
We first expand the characters and geometric products in the tail
variables \(y\), and then expand their rational coefficients in \(q\)
as Laurent series at \(q=0\). In particular, \(D(q)^{-1}\) has only
nonnegative powers.

\subsection{Coefficient extraction and truncation}
The normalized alternating Schur formula gives
\begin{equation}\label{eq:bandextract}
 Q(\lambda)-F(\lambda)
 =[q^By^\eta]\mathcal C(q,y)\mathcal A(y)
                       (\mathcal Q_b-\mathcal F_b)(1,q,y).
\end{equation}
Indeed in six variables it extracts the coefficient of \(x^\lambda\)
after multiplication by \(\prod_{i<j}(1-x_j/x_i)\). Distinct decreasing
shifted partitions cannot be permutations of each other, so only the
Schur term at \(\lambda\) contributes. All Laurent monomials have total
degree \(6b\); setting the first variable to one loses no information
about its exponent once the remaining five are fixed.

Give each \(y_\ell\) degree one and \(q\) degree zero.
The monomials of \(\mathcal A\) have total \(y\)-degree zero; those of
\(\mathcal C\) have nonnegative total \(y\)-degree. Character terms above
degree \(27\) therefore cannot contribute to \eqref{eq:bandextract}.
This uses total degree even when individual exponents in the alternant
are negative. Write \(\equiv_{27}\) for equality of coefficients of
nonnegative \(y\)-monomials through total degree \(27\), with coefficients
in \(\mathbb Q(q)\).

\subsection{Seven numerators independent of the degree}
For \(0\le j\le6\) define
\[
 L_j(q,y)=
 \prod_{\substack{\tau\in\NN^4,\ 1\le|\tau|\le6\\0\le u\le6-|\tau|}}
                         (1-y^\tau q^{u-j})^{-1}.
\]
Then
\begin{equation}\label{eq:targetband}
 \mathcal Q_b(1,q,y)\equiv_{27}
 \frac1{D(q)}\sum_{j=0}^6q^{jb}(-1)^j
                  e_j(q,q^2,\ldots,q^6)L_j(q,y).
\end{equation}
To prove this, fix a multiset of \(k\) inner monomials with nonzero tail.
The remaining \(b-k\) entries come from \(1,q,\ldots,q^6\), with
generating polynomial
\[
 h_{b-k}(1,q,\ldots,q^6)=
                    \prod_{h=1}^6\frac{1-q^{b-k+h}}{1-q^h}
                                      \quad(b-k\ge0).
\]
This Gaussian polynomial counts partitions in a \(6\)-by-\((b-k)\)
rectangle. In general their generating polynomial \(P_{r,s}\), with at
most \(r\) rows of length at most \(s\), satisfies
\(P_{r,s}=P_{r,s-1}+q^sP_{r-1,s}\), by removing a row of length \(s\)
when one is present, and \(P_{0,s}=P_{r,0}=1\).
The displayed product with \(r=b-k,s=6\) satisfies the same recurrence
and boundary values, proving the formula.
Expanding the numerator gives the \(e_j\) in \eqref{eq:targetband};
the factor \(q^{-jk}\) contributes \(q^{-j}\) per selected tail monomial,
and hence produces \(L_j\).
Within the cutoff \(k\le27\) and \(b\ge26\), so the only possible
negative value of \(b-k\) is \(-1\). Its numerator has the factor
\(1-q^0\) and vanishes, exactly as required. Thus the formula includes
the boundary \(b=26,k=27\); its cancellation is across the seven
\(j\)-terms.

Similarly, the binary complement of one tail monomial gives
\begin{equation}\label{eq:binarysource}
 h_b(1,q,y)\equiv_{27}
 \frac{\prod_\ell(1-y_\ell)^{-1}
       -q^{b+1}\prod_\ell(1-y_\ell/q)^{-1}}{1-q}.
\end{equation}
At tail degree \(t\), its coefficient is
\((1-q^{b+1-t})/(1-q)\). It counts the binary monomials when \(t\le b\)
and is zero at \(t=b+1\), the only extra possible degree.

For a cycle type \(I=(i_1,\ldots,i_v)\vdash6\), let
\[
 n_I=\frac{720}{\prod_{r\ge1}r^{m_r}m_r!},\qquad
 E_I(q)=\frac{D(q)}{\prod_{h=1}^v(1-q^{i_h})},
\]
where \(m_r\) counts cycles of length \(r\).
The quotient \(E_I\) is an integer polynomial of degree \(15\):
the multiplicity of the cyclotomic factor \(\Phi_d\) in its denominator
is the number of cycle lengths divisible by \(d\), which is at most
\(\lfloor6/d\rfloor\), its multiplicity in \(D\).
Apply \eqref{eq:binarysource} with variables raised to each cycle length
in \eqref{eq:cycle}. Together with \eqref{eq:targetband} this yields
\begin{equation}\label{eq:seven}
 720D(q)(\mathcal Q_b-\mathcal F_b)(1,q,y)
                         \equiv_{27}\sum_{j=0}^6q^{jb}G_j(q,y),
\end{equation}
where the seven numerators are explicitly
\begin{align}
G_j(q,y)={}&720(-1)^je_j(q,\ldots,q^6)L_j(q,y)\notag\\
 &-\sum_{I\vdash6}n_I
   \sum_{\substack{J\subseteq\{1,\ldots,v\}\\\sum_{h\in J}i_h=j}}
   (-1)^{|J|}q^j E_I(q)
   \prod_{\ell=1}^4\prod_{h=1}^v
       (1-y_\ell^{i_h}q^{-i_h\one_{h\in J}})^{-1}.
                                                     \label{eq:G}
\end{align}
Here \(\one_{h\in J}\) is the indicator of membership.
Every selected second term of \eqref{eq:binarysource} contributes
\(-q^{i_h(b+1)}\), explaining the sign and the residual \(q^j\)
after \(q^{jb}\) is factored out. Equal-length cycles still have distinct
positions in the subset sum. These formulas involve eleven cycle types
and seven values of \(j\), regardless of \(b\).

\subsection{Computing the numerators}
Write \(L_{j,a}=[y^a]L_j\). Euler differentiation of the product gives
\begin{equation}\label{eq:euler}
 L_{j,0}=1,\qquad
 |a|L_{j,a}=
 \sum_{\tau,u}\ \sum_{\substack{i\ge1\\i\tau\le a}}
              |\tau|q^{i(u-j)}L_{j,a-i\tau}\quad(|a|>0).
\end{equation}
For example, this follows by differentiating the logarithmic series
\(\sum_{i\ge1}(y^\tau q^{u-j})^i/i\). Its denominator \(i\) cancels
the extra \(i\) from differentiation. Every lookup has smaller total
tail degree, and the division by \(|a|\) is exact.

For the source product in \eqref{eq:G}, put
\[
 A_J(m)=[z^m]\prod_{h\in J}(1-z^{i_h})^{-1}.
\]
At a single tail-coordinate exponent \(m\), its Laurent coefficient is
\[
          \sum_{p=0}^m A_J(p)A_{\bar J}(m-p)q^{-p}.
\]
Multiply these four one-coordinate polynomials and then
\((-1)^{|J|}n_Iq^jE_I(q)\). These are the \texttt{ways} arrays and
convolutions of \texttt{source} in Appendix~\ref{app:band}.
The routines \texttt{target} and \texttt{gen\_source} implement
\eqref{eq:euler} and the cycle-type enumeration.

All raw arrays are symmetric in the four tail variables, so they need
only be stored at decreasing nonnegative exponent vectors; arbitrary
orders are looked up by sorting, without a sign.
To multiply by the cross factors in \(\mathcal C\), the coefficient
at a shift \(v\in\{0,1,2\}^4\) is
\[
                  (-1)^{n_1}q^{-n_2}(1+q^{-1})^{n_1},
       \qquad n_s=\#\{\ell:v_\ell=s\}.
\]
An in-place sweep in decreasing tail degree reads unmodified lower-degree
rows. A subsequent decreasing-exponent sweep multiplies by \(1-q\).
This is the routine \texttt{cross}.

The cross factors preserve symmetry in the four tail variables. To
extract a Schur coefficient, we now apply the alternating factor
\(\mathcal A\); the permutation signs enter at this step, rather than in the
sorted lookups.
Put \(\rho=(3,2,1,0)\), and let \(M_{j,a}=[y^a]\mathcal C G_j\),
with value zero at a vector with a negative coordinate.
The exact row produced by \texttt{schur\_row} is
\begin{equation}\label{eq:row}
 R_{j,\eta}(q)=\frac1{D(q)}
       \sum_{\pi\in\SSS_4}\sgn(\pi)
                     M_{j,\eta+\rho-\pi\rho}(q).
\end{equation}
This follows from
\(\mathcal A(y)=y^{-\rho}\sum_\pi\sgn(\pi)y^{\pi\rho}\).
Every accepted index in \eqref{eq:row} has nonnegative entries and total
degree \(|\eta|\le27\); in particular every coordinate is at most \(27\).
There are \(1908\) decreasing four-coordinate tails in this range.
Equations~\eqref{eq:bandextract} and \eqref{eq:seven} show that the final
integer tested is precisely
\begin{equation}\label{eq:finalinteger}
 \sum_{j=0}^6[q^{B-jb}]R_{j,\eta}(q)
       =720\bigl(Q(6b-B-|\eta|,B,\eta)-F(6b-B-|\eta|,B,\eta)\bigr).
\end{equation}
The notation on the right specifies partitions rather than gap vectors.

\subsection{The Laurent window}
At total tail degree \(t\), the coefficients and individual Euler
summands of \(L_j\) have \(q\)-support in \([-6t,5t]\).
The elementary factor adds at most \(21\). A cross term consuming
\(s\) tail units shifts down by at most \(s\), while its raw input has
degree \(t-s\); the factor \(1-q\) adds at most one.
Thus all target numerator summands lie in
\[
                      [-6t,5t+22].
\]
The source has raw support in \([-t,21]\) and final support in
\([-t,22]\), using \(\deg E_I=15\) and \(0\le j\le6\).
Alternation does not change total tail degree or \(q\)-support.
At \(t\le27\) all these supports are therefore contained in
\[
                       [-162,157]\subset[-168,161],
\]
the interval stored by the program. Its clipping operation discards
only zero terms. In particular all terms of the boundary cancellation
at \(b=26,k=27\) are retained separately.

Since \(D(0)=1\), its inverse is a power series with no negative
exponents. Starting from exponent \(-168\), six forward passes
\(a_n\gets a_n+a_{n-h}\), for \(h=1,\ldots,6\), divide by \(D\)
exactly. The numerator is zero above \(161\).
Expansion through \(447\) suffices because
\[
                         B-jb\le B\le3b\le447.
\]
Entries below \(-168\) are zero; every accessed nonnegative array index
is at most \(168+447=615\), within the \(616\) allocated entries.
Neither division nor the shift \(q^{jb}\) can bring a discarded high
exponent down to a requested one. The program converts to
arbitrary-precision integers before the denominator passes and uses them
for all final sums and divisibility tests.

Finally, \texttt{run\_thin} visits each of the \(1908\) tails once,
constructs its seven rows, and loops over the displayed ranges of
\(b\) and \(B\). For a fixed \(b,\eta\), the number of permitted
second parts is
\[
       3b-\left\lceil\frac{|\eta|}{2}\right\rceil-\eta_1+1.
\]
This is positive throughout the band. Summing it over the tails and
\(26\le b\le149\) gives \(57065668\) distinct partitions; there is
no dependence on which partitions the first program flagged.

\begin{verification}\label{ver:band}
The complete program in Appendix~\ref{app:band} checks \(57065668\)
coefficients in this band. Of these, \(177134\) are zero and none are
negative; every integer in \eqref{eq:finalinteger} is divisible by \(720\).
The grouped counts appear in Table~\ref{tab:band}.
\end{verification}

\section{Finite arithmetic, reproduction, and conclusion}\label{sec:verification}

\subsection{Indexing and arithmetic in the certificate program}
Appendix~\ref{app:certificate} lists the complete program used for
Verification~\ref{ver:certificates}; it has no external data input.
Its partition order reads parts from last to first. If \(p_j(u)\)
counts partitions of \(u\) into at most \(j\) parts, the prefix table
counts those whose last padded part is \(<k\), by
\[
 N_j(u,k)=N_j(u,k-1)+p_{j-1}(u-j(k-1)).
\]
The second summand subtracts the last part from all \(j\) positions.
Negative arguments count as zero and \(p_0(0)=1\).
Subtracting two prefix counts ranks the next part after the already
fixed smaller neighbor. A shorter partition padded by zeros has the
same rank. For the tail-restricted table, first group by its exact last
three parts, subtract their largest part from the preceding three, and
use the three-part rank. The grid similarly groups by \(w(m)\) and ranks
the partition corresponding to its first five gaps.

The allocated prefix array has size parameter \(1250\). Its terminal
weight need not supply a complete unrestricted count; all counts used
for multiplicities and grid indices have weight at most \(894\), strictly
below that terminal boundary. The prefix entries are bounded by
\(\binom{1254}{5}<2^{63}\).
The actual largest complete levels used in six, four, and three variables
have respectively \(1229120\), \(261072\), and \(67051\) partitions.
The largest tail-restricted level has \(1263138\) entries, and the grid
has \(6611697\). These ranks, offsets, and predecessor indices fit signed
32-bit integers. Independently of these evaluated counts, coarse bounds
for the three complete levels, the grid, and any tail level are
\[
 \max\left\{\binom{155}{5},\binom{333}{3},\binom{896}{2},
             \binom{154}{5},\binom{21}{3}\binom{272}{2}\right\}<2^{31}.
\]
These count unrestricted weak compositions. For the grid,
\(w(m)\le149\) implies \(\sum_i m_i\le149\); for a tail level there
are at most \(\binom{21}{3}\) ordered tails of size at most \(18\), and
each remaining three-part prefix has weight at most \(270\).

Multiplicity arithmetic is signed 128-bit. The target recurrences check
every accumulation for overflow, and test integrality and nonnegativity
before division. For the source, only \(b\le25\) is needed.
Each intermediate multiplicity is at most \(151^{10}\): it is bounded
by its multiplicity in \((\Sym^b V)^{\otimes j}\), \(j\le6\), and the
successive interlacing partitions have at most
\(151^{0+1+2+3+4}\) choices. In one source Newton step, there are fewer
than \(6\cdot151^4+6\) strip terms, by fixing one coordinate with the
size constraint. Thus every signed accumulation is bounded absolutely by
\[
                   151^{10}(6\cdot151^4+6)<2^{127}.
\]
For \(U\), every width in every actual or high-residue call is at most
\(920\). A chain has at most \(1+2+3+4=10\) width factors, so
\[
                      B_I\le920^{10}<Z=2^{100}.
\]
The minimum initialized at \(Z\) therefore never clips an upper bound,
and summing class weights costs a factor at most \(720\).

For the power bound, \(S\) has
\(\dim\Sym^6\CC^6=\binom{11}{5}=462\) polynomial generators.
Consequently \(\binom{d_0+461}{461}\ge Q(x)\), so the binomial search
for \(z\), before its cap at \(18\), stops by \(461\).
With \(C\le Z\) and \(d_0\le149\), its first crossing is at most
\(150Z\); predivision multiplication costs at most another factor \(610\),
still below \(2^{127}\). After the cap, the other binomials are at most
\(\binom{167}{18}\). Products in the final lower bound are screened
against \(Z\) before multiplication. Capping can only weaken a lower
certificate. These observations cover the unchecked arithmetic as well
as the program's explicit assertions.

OpenMP distributes distinct coefficient destinations within one degree.
All previous degrees are read-only, and each work-sharing loop has its
barrier before the next degree. The mathematical result is therefore
independent of the thread count.

\subsection{Integer bounds for the band program}
It remains to bound the signed 128-bit arithmetic before the
arbitrary-precision division described in Section~\ref{sec:band}.
For a Laurent series truncated to total tail degree at most \(27\),
use the sum of absolute coefficients as its norm.
The norm of \(L_j\) is bounded by
\[
 N=\sum_{n=0}^{27}[z^n]P(z),\qquad
 P(z)=\prod_{h=1}^6(1-z^h)^{-M_h},\qquad
 M_h=(7-h)\binom{h+3}{3}.
\]
Here \((M_1,\ldots,M_6)=(24,50,80,105,112,84)\).
There is a convenient purely elementary estimate
\begin{equation}\label{eq:normbound}
                         N\le4^{27}P(1/4)<2^{79}.
\end{equation}
The first inequality follows from nonnegative coefficients.
For the second, put \(r_h=4^h/2\). Bernoulli's inequality gives
\((1-4^{-h})^{r_h}>1/2\).
Thus the \(h\)-th factor of \(P(1/4)\) is at most
\(2^{\lceil M_h/r_h\rceil}\), with strict total inequality.
The six ceilings are \(12,7,3,1,1,1\), whose sum is \(25\);
combine this with \(4^{27}=2^{54}\).

Each four-coordinate source product has at most \(24\) geometric factors,
each consuming at least one tail unit, so its norm is at most
\[
                         S_0=\binom{51}{24}<2^{51}.
\]
Using \(\|D\|_1\le64\), the quotient \(E_I\) has norm at most
\[
             Q_0=64\binom{21}{6}=3472896<2^{22}.
\]
Indeed dominate all reciprocal cycle factors by \((1-q)^{-6}\) and
sum through degree \(15\). The partial quotient calculations through
degree \(21\) have bound \(64\binom{27}{6}=18944640\), and the
one-coordinate arrays \(A_J\) have coefficients at most
\(\binom{32}{5}=201376\). Their additions fit signed 32-bit integers.
Class-size divisions are exact: for each partial cycle multiset of
size \(r\le6\), its centralizer order divides \(r!\), hence divides \(720\).

The undivided Euler step costs at most \(27N\).
The absolute weight of all target elementary masks is \(720\cdot64\),
and the class-subset weight is
\[
 \sum_{I\vdash6}n_I2^{v(I)}
  =\sum_{\sigma\in\SSS_6}2^{\#\text{cycles}(\sigma)}
  =2\cdot3\cdot4\cdot5\cdot6\cdot7=5040.
\]
The last identity follows inductively by inserting the new largest
element into an existing cycle or making it a new cycle.
Cross multiplication including \(1-q\) has norm at most \(512\);
alternation costs at most \(24\). A bound for every predivision
coefficient, partial sum, and product is therefore
\begin{equation}\label{eq:allbounds}
        512\cdot24\bigl(27\cdot720\cdot64N+5040Q_0S_0\bigr)
                         <2^{115}<2^{127}.
\end{equation}
The source convolutions have nonnegative factors before their final
signs, so this also bounds their intermediate products.
Potentially large products in the code already have a 128-bit operand.
The row divisions and final sums then use arbitrary precision.
Each output counter covers at most six degrees, \(1908\) tails, and
\(448\) second parts per degree and tail. Thus even the coarse bound
\(1908\cdot6\cdot448<2^{31}\) keeps these counters within signed
32-bit range.

\subsection{Recorded computations and reproduction}
The reference output streams contain \(20\) base rows, \(21\) flag rows,
and \(21\) band rows. Tables~\ref{tab:base} and~\ref{tab:band} present
all \(62\) numeric rows, including every tested degree. The reproduction
script checks these streams against fixed hashes and compares their
numeric contents with both manuscript tables.

A full reproduction compiled the two printed sources with
GNU C++~13.3.0 and executed both programs over their entire stated ranges,
with assertions enabled. The first used two OpenMP threads and the
second was serial. Both exited with status zero, and all four output
streams matched the reference streams byte for byte. This execution
reproduced all \(62\) numeric rows in the tables.

The programs retain the complete ranges \(6\le b\le149\).
For the main proof using Corollary~\ref{cor:quadratic-range}, only the
base comparisons \(6\le b\le25\) and the first row of
Table~\ref{tab:band}, for \(26\le b\le29\), are required.
That row contains \(29814\) flags and \(467068\) band coefficients;
the band differences are nonnegative, with \(2114\) equalities.
The later rows give the additional finite checks for the independent
route through Proposition~\ref{prop:large}. Exact target tables at higher
degrees are computed from the recurrences, not assumed from Foulkes'
conjecture, so their use in lower certificates introduces no circularity.

\begin{table}[ht]
\centering
\begin{tabular}{r|r|r|r|r}
\(b\) range&flagged pairs&largest \(t_2\)&band coefficients&zero differences\\ \hline
26--29&29814&25&467068&2114\\
30--35&46805&25&872322&3621\\
36--41&50091&25&1078386&4161\\
42--47&54443&25&1284450&4701\\
48--53&63027&25&1490514&5241\\
54--59&74729&20&1696578&5781\\
60--65&90684&18&1902642&6321\\
66--71&109505&18&2108706&6861\\
72--77&130920&19&2314770&7401\\
78--83&155274&20&2520834&7941\\
84--89&182871&21&2726898&8481\\
90--95&214444&22&2932962&9021\\
96--101&248073&23&3139026&9561\\
102--107&285382&24&3345090&10101\\
108--113&324614&25&3551154&10641\\
114--119&366623&26&3757218&11181\\
120--125&410723&26&3963282&11721\\
126--131&457715&26&4169346&12261\\
132--137&506088&26&4375410&12801\\
138--143&557588&26&4581474&13341\\
144--149&610285&27&4787538&13881\\ \hline
total&4969698&27&57065668&177134
\end{tabular}
\caption{Full retained verification range. The first row is the residual
range required with quadratic stabilization; subsequent rows support the
independent chart route. All flags lie in the checked band, and no
negative difference is found. The program's first row is labeled \(24\), the
multiple of six below the first tested degree; the explicit degree
ranges here avoid that convention.}\label{tab:band}
\end{table}

The supporting package contains the complete program sources, reference
output streams, their hashes and numeric tables, and a reproduction
script. A new full run creates its own compiler and execution metadata.
The commands, run from the paper directory, are
\begin{verbatim}
python3 -B verification/verify_computations.py --check
python3 -B verification/verify_computations.py --run --threads 2
\end{verbatim}
The first checks source and reference-stream identities, all \(62\)
numeric rows, agreement with the manuscript tables, and the numerical
bounds above; it does not execute the mathematical programs.
The second compiles both unchanged sources with
\texttt{-std=c++17 -O3 -UNDEBUG}, enables OpenMP only for the first,
executes them, and requires both complete output streams to match
the reference streams byte for byte. It keeps each run in a new external
directory, with the compiler version, build and execution commands,
exit statuses, timings, hashes, and captured streams. The option
\texttt{--serial} may replace \texttt{--threads 2} to compile and run
the first program without OpenMP.
The two C++ files are also printed in full below, so every finite step
is specified in this paper. The coefficient formulas, indexing
arguments, and arithmetic bounds establish what they check; successful
execution supplies the finite nonnegativity assertions.

\subsection{Completion of the proof}
\begin{proof}[Proof of Theorem~\ref{thm:main}]
Corollary~\ref{cor:quadratic-range} settles \(b\ge30\).
Verification~\ref{ver:certificates} settles \(6\le b\le25\).
Proceed by increasing \(b\) from \(26\) through \(29\).
If \(d_6\ge6\), Lemma~\ref{lem:rectangles} reduces the comparison to
degree \(b-6\ge20\). Otherwise \(d=r+6m\) lies in the exhaustive
classification of Section~\ref{sec:certificates}. Every pair is either
cleared there by a proved certificate or the chart, or flagged.
The flags have \(t_2\le27\) and are covered by
Verification~\ref{ver:band}. Hence \(Q(\lambda)\ge F(\lambda)\) for
every source partition and every \(b\ge6\).
The source has no constituents of length greater than six;
specialization at extra variables zero transfers the comparisons to
every finite dimension. Complete reducibility gives the required
equivariant injection.
\end{proof}

\begin{remark}[The independent full-range route]
The same finite argument, with no change of predicate or arithmetic,
continues from \(b=30\) through \(b=149\), as recorded in the full
tables. Combining it with Proposition~\ref{prop:large} for \(b\ge150\)
proves the theorem without using the quadratic companion.
The chart and certificate methods are therefore retained as a complete
independent proof, while Corollary~\ref{cor:quadratic-range} gives the
shorter large-degree route.
\end{remark}

\bibliographystyle{plain}
\bibliography{references}
\appendix
\section{The bounded certificate program}\label{app:certificate}
The following is the complete first program. It uses GNU C++17 signed
128-bit integers and optional OpenMP; assertions must remain enabled.
The source file is \nolinkurl{verification/programs/appendix-a.cpp}.
See Sections~\ref{sec:upper}--\ref{sec:certificates} for the construction
and Section~\ref{sec:verification} for arithmetic and reproduction.
\lstinputlisting[style=proofcode]{../verification/programs/appendix-a.cpp}
\section{The tail-band program}\label{app:band}
The following is the complete second program. It uses GNU C++17 and
Boost.Multiprecision for the final coefficient sums.
The source file is \texttt{verification/programs/appendix-b.cpp}.
See Section~\ref{sec:band} for the construction and
Section~\ref{sec:verification} for arithmetic and reproduction.
\lstinputlisting[style=proofcode]{../verification/programs/appendix-b.cpp}
\end{document}
